% Question 1.(b) -- Lagrange function, dual problem and KKT conditions of the hourly problem
%
% SCAFFOLD ONLY: structure and \todo prompts, no derivation. Fill each \todo with your own work.
% Task text (course_materials/Assignment_1_Instructions.html, Question 1.(b)):
%   "Formulate the Lagrange function, the dual problem, and the KKT conditions of the hourly
%    problem. Does strong duality hold? Are the KKT conditions necessary and/or sufficient
%    optimality conditions? Justify your answer from the properties stated in (a)."
% Grading (Table 2, row 1.(b)): mathematical model + writing; no code.
% Source material: Lecture 3 notes -- sections 3 (Lagrangian), 4 (Lagrange duality),
%   5 (LP duality), 6 (KKT).
%
% Called from the end of sources/1_.tex with % Question 1.(b) -- Lagrange function, dual problem and KKT conditions of the hourly problem
%
% SCAFFOLD ONLY: structure and \todo prompts, no derivation. Fill each \todo with your own work.
% Task text (course_materials/Assignment_1_Instructions.html, Question 1.(b)):
%   "Formulate the Lagrange function, the dual problem, and the KKT conditions of the hourly
%    problem. Does strong duality hold? Are the KKT conditions necessary and/or sufficient
%    optimality conditions? Justify your answer from the properties stated in (a)."
% Grading (Table 2, row 1.(b)): mathematical model + writing; no code.
% Source material: Lecture 3 notes -- sections 3 (Lagrangian), 4 (Lagrange duality),
%   5 (LP duality), 6 (KKT).
%
% Called from the end of sources/1_.tex with % Question 1.(b) -- Lagrange function, dual problem and KKT conditions of the hourly problem
%
% SCAFFOLD ONLY: structure and \todo prompts, no derivation. Fill each \todo with your own work.
% Task text (course_materials/Assignment_1_Instructions.html, Question 1.(b)):
%   "Formulate the Lagrange function, the dual problem, and the KKT conditions of the hourly
%    problem. Does strong duality hold? Are the KKT conditions necessary and/or sufficient
%    optimality conditions? Justify your answer from the properties stated in (a)."
% Grading (Table 2, row 1.(b)): mathematical model + writing; no code.
% Source material: Lecture 3 notes -- sections 3 (Lagrangian), 4 (Lagrange duality),
%   5 (LP duality), 6 (KKT).
%
% Called from the end of sources/1_.tex with % Question 1.(b) -- Lagrange function, dual problem and KKT conditions of the hourly problem
%
% SCAFFOLD ONLY: structure and \todo prompts, no derivation. Fill each \todo with your own work.
% Task text (course_materials/Assignment_1_Instructions.html, Question 1.(b)):
%   "Formulate the Lagrange function, the dual problem, and the KKT conditions of the hourly
%    problem. Does strong duality hold? Are the KKT conditions necessary and/or sufficient
%    optimality conditions? Justify your answer from the properties stated in (a)."
% Grading (Table 2, row 1.(b)): mathematical model + writing; no code.
% Source material: Lecture 3 notes -- sections 3 (Lagrangian), 4 (Lagrange duality),
%   5 (LP duality), 6 (KKT).
%
% Called from the end of sources/1_.tex with \input{sources/1b_duality}.
% Whatever multipliers you introduce: add them to sources/9_nomenclat.tex (with units) and, if
% you create macros for them, put the \newcommand next to the others in main.tex.

\subsection{(b) Lagrangian, dual problem and KKT conditions}

\subsubsection{Setting up}
\paragraph{Hourly problem in the form used here.}
\todo{Restate the hourly problem of (a) for a generic hour $t$ (one line, refer to the
equations of (a)). Decide and state your conventions: do you keep it as a maximisation or
turn it into a minimisation as in the lecture notes? How is each constraint written
($\leq 0$, $\geq 0$, $= 0$)? The signs of all multipliers below depend on these choices, so
fix them once here.}

\paragraph{Multipliers.}
\todo{Associate one dual variable to each constraint: give a table or list (name, the
constraint it belongs to, its sign restriction, its unit). Decide which constraints get a
multiplier and which ones stay as variable bounds ($x \geq 0$) -- say which and why.}

\subsubsection{Lagrange function}
\todo{Write the Lagrangian of the hourly problem with your multipliers.}

\subsubsection{Dual problem}
\paragraph{Dual function.}
\todo{Write the dual function (minimise/maximise the Lagrangian over the primal variables).
State the condition on the multipliers for which it is finite, and where it comes from.}

\paragraph{Dual problem.}
\todo{Write the dual problem: objective, constraints on the multipliers. Which problem class
is it?}

\subsubsection{KKT conditions}
\todo{Write the four groups of conditions for your problem, one line each per variable or
constraint: (i) stationarity, (ii) primal feasibility, (iii) dual feasibility, (iv)
complementary slackness.}

\subsubsection{Strong duality}
\todo{Does strong duality hold? Justify from the properties of the problem stated in (a)
(problem class, feasibility, boundedness), and state what it implies for the primal and dual
optimal values.}

\subsubsection{Necessity and sufficiency of the KKT conditions}
\todo{Are the KKT conditions necessary, sufficient, both, or neither for optimality of this
problem? Justify from the properties stated in (a), separately for necessity and for
sufficiency.}
.
% Whatever multipliers you introduce: add them to sources/9_nomenclat.tex (with units) and, if
% you create macros for them, put the \newcommand next to the others in main.tex.

\subsection{(b) Lagrangian, dual problem and KKT conditions}

\subsubsection{Setting up}
\paragraph{Hourly problem in the form used here.}
\todo{Restate the hourly problem of (a) for a generic hour $t$ (one line, refer to the
equations of (a)). Decide and state your conventions: do you keep it as a maximisation or
turn it into a minimisation as in the lecture notes? How is each constraint written
($\leq 0$, $\geq 0$, $= 0$)? The signs of all multipliers below depend on these choices, so
fix them once here.}

\paragraph{Multipliers.}
\todo{Associate one dual variable to each constraint: give a table or list (name, the
constraint it belongs to, its sign restriction, its unit). Decide which constraints get a
multiplier and which ones stay as variable bounds ($x \geq 0$) -- say which and why.}

\subsubsection{Lagrange function}
\todo{Write the Lagrangian of the hourly problem with your multipliers.}

\subsubsection{Dual problem}
\paragraph{Dual function.}
\todo{Write the dual function (minimise/maximise the Lagrangian over the primal variables).
State the condition on the multipliers for which it is finite, and where it comes from.}

\paragraph{Dual problem.}
\todo{Write the dual problem: objective, constraints on the multipliers. Which problem class
is it?}

\subsubsection{KKT conditions}
\todo{Write the four groups of conditions for your problem, one line each per variable or
constraint: (i) stationarity, (ii) primal feasibility, (iii) dual feasibility, (iv)
complementary slackness.}

\subsubsection{Strong duality}
\todo{Does strong duality hold? Justify from the properties of the problem stated in (a)
(problem class, feasibility, boundedness), and state what it implies for the primal and dual
optimal values.}

\subsubsection{Necessity and sufficiency of the KKT conditions}
\todo{Are the KKT conditions necessary, sufficient, both, or neither for optimality of this
problem? Justify from the properties stated in (a), separately for necessity and for
sufficiency.}
.
% Whatever multipliers you introduce: add them to sources/9_nomenclat.tex (with units) and, if
% you create macros for them, put the \newcommand next to the others in main.tex.

\subsection{(b) Lagrangian, dual problem and KKT conditions}

\subsubsection{Setting up}
\paragraph{Hourly problem in the form used here.}
\todo{Restate the hourly problem of (a) for a generic hour $t$ (one line, refer to the
equations of (a)). Decide and state your conventions: do you keep it as a maximisation or
turn it into a minimisation as in the lecture notes? How is each constraint written
($\leq 0$, $\geq 0$, $= 0$)? The signs of all multipliers below depend on these choices, so
fix them once here.}

\paragraph{Multipliers.}
\todo{Associate one dual variable to each constraint: give a table or list (name, the
constraint it belongs to, its sign restriction, its unit). Decide which constraints get a
multiplier and which ones stay as variable bounds ($x \geq 0$) -- say which and why.}

\subsubsection{Lagrange function}
\todo{Write the Lagrangian of the hourly problem with your multipliers.}

\subsubsection{Dual problem}
\paragraph{Dual function.}
\todo{Write the dual function (minimise/maximise the Lagrangian over the primal variables).
State the condition on the multipliers for which it is finite, and where it comes from.}

\paragraph{Dual problem.}
\todo{Write the dual problem: objective, constraints on the multipliers. Which problem class
is it?}

\subsubsection{KKT conditions}
\todo{Write the four groups of conditions for your problem, one line each per variable or
constraint: (i) stationarity, (ii) primal feasibility, (iii) dual feasibility, (iv)
complementary slackness.}

\subsubsection{Strong duality}
\todo{Does strong duality hold? Justify from the properties of the problem stated in (a)
(problem class, feasibility, boundedness), and state what it implies for the primal and dual
optimal values.}

\subsubsection{Necessity and sufficiency of the KKT conditions}
\todo{Are the KKT conditions necessary, sufficient, both, or neither for optimality of this
problem? Justify from the properties stated in (a), separately for necessity and for
sufficiency.}
.
% Whatever multipliers you introduce: add them to sources/9_nomenclat.tex (with units) and, if
% you create macros for them, put the \newcommand next to the others in main.tex.

\subsection{(b) Lagrangian, dual problem and KKT conditions}

\subsubsection{Setting up}
\paragraph{Hourly problem in the form used here.}
\todo{Restate the hourly problem of (a) for a generic hour $t$ (one line, refer to the
equations of (a)). Decide and state your conventions: do you keep it as a maximisation or
turn it into a minimisation as in the lecture notes? How is each constraint written
($\leq 0$, $\geq 0$, $= 0$)? The signs of all multipliers below depend on these choices, so
fix them once here.}

\paragraph{Multipliers.}
\todo{Associate one dual variable to each constraint: give a table or list (name, the
constraint it belongs to, its sign restriction, its unit). Decide which constraints get a
multiplier and which ones stay as variable bounds ($x \geq 0$) -- say which and why.}

\subsubsection{Lagrange function}
\todo{Write the Lagrangian of the hourly problem with your multipliers.}

\subsubsection{Dual problem}
\paragraph{Dual function.}
\todo{Write the dual function (minimise/maximise the Lagrangian over the primal variables).
State the condition on the multipliers for which it is finite, and where it comes from.}

\paragraph{Dual problem.}
\todo{Write the dual problem: objective, constraints on the multipliers. Which problem class
is it?}

\subsubsection{KKT conditions}
\todo{Write the four groups of conditions for your problem, one line each per variable or
constraint: (i) stationarity, (ii) primal feasibility, (iii) dual feasibility, (iv)
complementary slackness.}

\subsubsection{Strong duality}
\todo{Does strong duality hold? Justify from the properties of the problem stated in (a)
(problem class, feasibility, boundedness), and state what it implies for the primal and dual
optimal values.}

\subsubsection{Necessity and sufficiency of the KKT conditions}
\todo{Are the KKT conditions necessary, sufficient, both, or neither for optimality of this
problem? Justify from the properties stated in (a), separately for necessity and for
sufficiency.}
